\section{Excluding lower-energy continuum states}
\label{sec:continuum}

We now compare the spectral infimum of the full Coulomb operator with
the half-filled Hubbard energy.  In particular, the comparison below
includes a lower bound on every state outside the selected orbital
space.

The proof has three distinct outputs. We first identify the one-particle
matrix in orthonormal site modes and prove a fixed positive gap on its
entire orthogonal complement. We then bound the many-electron coupling
to that complement and eliminate it by a quadratic-form square completion.
Only after this all-state lower bound is available do we compare the
compressed Coulomb interaction with the Hubbard onsite term. Keeping
these steps separate prevents a variational upper bound from being used
as a NO-instance lower bound.

We use the rational centers and tuned parameters constructed in
Sections~\ref{sec:wells} and~\ref{sec:hopping}.  Increasing the fixed
polynomial choice of \(D\) if necessary, we may assume throughout this
section that
\[
 m+S\le D^{1/4},\qquad D\ge D_0,
\]
where the constant \(D_0\) can be chosen uniformly.  The estimates from
Propositions~\ref{prop:isolated-well} and~\ref{prop:tuning} and
Lemma~\ref{lem:well-tail} that we will use are
\begin{align}
 \|\psi_i-\phi_i\|_{H^2}&\le C m/D,
 &\|\psi_i\|_{H^2}&\le C, \label{eq:continuum-orbital-input}\\
 |\lambda_i+1/2|&\le e^{-1.01D},
 &|\psi_i(y)|&\le C e^{-.99(r_i-.01D)}. \label{eq:continuum-tail-input}
\end{align}
There is also a constant \(g_{\mathrm w}>0\), independent of the
construction parameters, such that
\begin{equation}
 h_i-\lambda_i\ge
 g_{\mathrm w}\bigl(I-|\psi_i\rangle\langle\psi_i|\bigr).
 \label{eq:continuum-well-gap}
\end{equation}
All operator inequalities involving an unbounded operator are
understood as quadratic-form inequalities.  By
Lemma~\ref{lem:contact-layout} and the displacement bound,
\begin{equation}
 |X_i-X_j|\ge D-2S\quad(i\ne j),\qquad
 |X_i-X_j|\ge\sqrt2D-2S\quad(\{i,j\}\notin\mathcal C).
 \label{eq:continuum-separations}
\end{equation}
The second inequality concerns distinct noncontact sites.

\subsection{The one-particle matrix}

Recall the hole potentials from \cref{sec:hopping}, and write
\[
 A=h+\tfrac12,\qquad
 V_i^{\mathrm{hole}}=(1-\chi_i)V_i,\qquad
 \epsilon_i=\lambda_i+\tfrac12.
\]
For the chosen orientation \(i\to j\) of a contact, define
\[
 B_{ij}=\langle\psi_j,V_j^{\mathrm{hole}}\psi_i\rangle,\qquad
 \epsilon_{\mathrm{hop}}
 =\max_{i\to j}\bigl|Z_*B_{ij}-t_{ij}\bigr|.
\]
The maximum is zero if there are no contacts.
Proposition~\ref{prop:calibration} gives
\begin{equation}
 \epsilon_{\mathrm{hop}}\le C\mathcal E_{\mathrm{cal}},
 \label{eq:continuum-calibration-interface}
\end{equation}
because all the target hoppings are at most one.
Let \(\mathsf T\) be the real symmetric \(m\)-by-\(m\) matrix with
entry \(-t_{ij}\) on contacts and zero elsewhere, including its
diagonal.

Define \(\mathcal U:\mathbb C^m\to L^2(\mathbb R^3)\) by
\(\mathcal U e_i=\psi_i\), and let \(G=\mathcal U^*\mathcal U\)
be its Gram matrix.  The notation \(\mathcal U\) denotes this map;
the scalar onsite repulsion remains \(U=5/8\).

\begin{lemma}[Residuals and the orthonormalized matrix]
\label{lem:one-body-matrix}
For the preceding parameters and sufficiently large \(D\),
\begin{equation}
 \|A\psi_i\|_2\le e^{-.75D},\qquad
 \|G-I\|\le e^{-.85D}.
 \label{eq:continuum-residual-gram}
\end{equation}
In particular \(G\) is positive definite.
The isometry
\(\mathcal J_0=\mathcal U G^{-1/2}\) has orthonormal columns
\(\omega_i=\mathcal J_0e_i\), spanning the same spatial orbital
space.  Its matrix satisfies
\begin{align}
 \|Z_*\mathcal J_0^*A\mathcal J_0-\mathsf T\|
 &\le \delta_1, \label{eq:continuum-one-body-error}\\
 \delta_1
 &:=m\epsilon_{\mathrm{hop}}+Z_*e^{-1.01D}
       +CmZ_*e^{-1.1D}+CmZ_*e^{-1.6D}. \notag
\end{align}
If \(P_0=\mathcal J_0\mathcal J_0^*\), then
\begin{equation}
 \|AP_0\|\le C\sqrt m\,e^{-.75D},
 \qquad
 \|\omega_i-\phi_i\|_{H^1}\le Cm/D.
 \label{eq:continuum-projected-residual}
\end{equation}
\end{lemma}

\begin{proof}
Since \(h=h_i-\sum_{k\ne i}V_k^{\mathrm{hole}}\),
\begin{equation}
 A\psi_i=\epsilon_i\psi_i-
       \sum_{k\ne i}V_k^{\mathrm{hole}}\psi_i.
 \label{eq:continuum-residual-identity}
\end{equation}
The \(k\)-th hole is supported in \(r_k\le D/8\).  On it,
\[
 r_i\ge 7D/8-2S\quad(k\ne i).
\]
Both Coulomb poles of \(V_k\) lie in this hole for large \(D\).
The charges \(q,q'\) are bounded, and direct radial integration,
enlarging the integration ball by at most four for the secondary
pole, gives
\begin{equation}
 \|V_k^{\mathrm{hole}}\|_2\le C D^{1/2},
 \qquad
 \|V_k^{\mathrm{hole}}\|_1\le C D^2.
 \label{eq:continuum-hole-norms}
\end{equation}
Consequently the sum in \eqref{eq:continuum-residual-identity}
has norm at most
\[
 CmD^{1/2}\exp(-.85635D+1.98S).
\]
Together with the tuning error, this is at most \(e^{-.75D}\)
for the stated regime and sufficiently large \(D\).

For \(i\ne j\), put \(\ell_{ij}=|X_i-X_j|\).
The product of the two tail bounds is at most
\[
 C\exp\bigl(-.99(r_i+r_j)+.0198D\bigr).
 \]
Use \(r_i+r_j\ge \ell_{ij}\) on a factor with exponent \(.90\);
the remaining factor is integrable uniformly in the centers,
since
\(\int e^{-.09(r_i+r_j)}\,\dd y\le
\int e^{-.09r_i}\,\dd y\le C\).
Thus
\[
 |G_{ij}|\le C e^{-.90\ell_{ij}+.0198D}
 \le C e^{-.8802D+1.8S}.
\]
The diagonal entries are one.  The row-sum estimate and
\(m+S\le D^{1/4}\) prove the second inequality in
\eqref{eq:continuum-residual-gram}.

Let \(M=\mathcal U^*A\mathcal U\), which is real symmetric.
We next estimate its entries before orthonormalization.
If \(k\ne i,j\), both \(r_i\) and \(r_j\) are at least
\(7D/8-2S\) on the \(k\)-th hole.  Equations
\eqref{eq:continuum-tail-input} and
\eqref{eq:continuum-hole-norms} therefore imply
\[
 \sum_{k\ne i,j}
 |\langle\psi_j,V_k^{\mathrm{hole}}\psi_i\rangle|
 \le CmD^2e^{-1.7127D+3.96S}.
\]
The same estimate bounds
\(\sum_{k\ne i}\langle\psi_i,V_k^{\mathrm{hole}}\psi_i\rangle\).
For a noncontact pair, the remaining hole at \(j\) is bounded using
\(r_i+r_j\ge\sqrt2D-2S\), giving
\[
 |\langle\psi_j,V_j^{\mathrm{hole}}\psi_i\rangle|
 \le CD^2
 e^{-(.99\sqrt2-.0198)D+1.98S}.
\]
Here \(.99\sqrt2-.0198>1.38\).
All these bounds are smaller than \(e^{-1.1D}\), after
absorbing their polynomial factors and fixed constants.
For \(i\ne j\), the term \(\epsilon_iG_{ji}\) is at most
\(e^{-1.86D}\).  It follows, adjusting a fixed constant, that
\begin{align}
 |M_{ii}-\epsilon_i|&\le Ce^{-1.1D},\notag\\
 |M_{ji}+B_{ij}|&\le Ce^{-1.1D}
       &&(i\to j\text{ a chosen contact}),\label{eq:continuum-raw-entries}\\
 |M_{ij}|&\le Ce^{-1.1D}
       &&(i\ne j,\ \{i,j\}\notin\mathcal C).\notag
\end{align}
The exactly symmetric matrix \(M\) requires only one of the two
orientations of each contact.  Its diagonal tuning contribution
has norm at most \(e^{-1.01D}\), not \(m e^{-1.01D}\).
Entrywise comparison with \(\mathsf T\) and a row-sum bound now give
\[
 \|Z_*M-\mathsf T\|
 \le m\epsilon_{\mathrm{hop}}
     +Z_*e^{-1.01D}+CmZ_*e^{-1.1D}.
\]

For completeness, the residual estimate controls the entire matrix,
without using a dimension-dependent bound on a many-body space:
\[
 \|A\mathcal U\|\le\sqrt m\,e^{-.75D},\qquad
 \|\mathcal U\|=\|G\|^{1/2}\le2,\qquad
 \|M\|\le C\sqrt m\,e^{-.75D}.
\]
Functional calculus on the positive matrix \(G\) gives
\(\|G^{-1/2}-I\|\le Ce^{-.85D}\) and
\(\|G^{-1/2}\|\le2\).  Consequently
\begin{equation}
 \|G^{-1/2}MG^{-1/2}-M\|
 \le C\sqrt m\,e^{-1.6D}
 \le Cm e^{-1.6D}.
 \label{eq:continuum-gram-correction}
\end{equation}
This orthonormalizes the matrix of the shifted operator
\(h+1/2\).  In the original nonorthogonal modes,
\(M=\mathcal U^*h\mathcal U+\tfrac12G\).
The shift is essential in this estimate: one must add
\(\tfrac12G\), rather than \(\tfrac12I\), before
orthonormalization.
Equation~\eqref{eq:continuum-one-body-error} follows.

The same residual estimate applied to
\(A\mathcal J_0=A\mathcal U G^{-1/2}\) proves the first part of
\eqref{eq:continuum-projected-residual}.
Finally, as a map into \(H^1\), \(\mathcal U\) has norm at most
\(C\sqrt m\), by \eqref{eq:continuum-orbital-input}.  Thus
\[
 \|\omega_i-\psi_i\|_{H^1}
 \le C\sqrt m\,e^{-.85D}.
\]
Combining this with \(\|\psi_i-\phi_i\|_{H^1}\le Cm/D\)
gives the remaining assertion.
\end{proof}

We now include spin.  Set
\[
 \mathfrak h=L^2(\mathbb R^3;\mathbb C^2),\qquad
 P=P_0\otimes I_2,\qquad
 \mathcal J=\mathcal J_0\otimes I_2.
\]
Thus \(P_0\) has rank \(m\), whereas \(P\) has rank \(2m\).
The operator \(A\) acts trivially on spin; we use the same symbol
for this extension.  Its residual and matrix bounds are unchanged.

\subsection{A gap on the entire one-particle complement}

Spatial localization separates the nuclei from the exterior, where the
attraction is uniformly weak. We use the IMS product-rule identity
(compare Simon~\cite[Lemma~3.1]{Simon1983}); its kinetic-error coefficient
is $1/2$ for our normalization $-\Delta/2$. The identity is derived below.
The additional point to verify is that a globally orthogonal function
has only exponentially small components along the local ground states
after localization.

\begin{proposition}[Continuum complement gap]
\label{prop:complement-gap}
There is a constant \(g>0\) such that, for all sufficiently large
\(D\),
\[
 (I-P)A(I-P)\ge g(I-P).
\]
The inequality holds on the full one-particle form domain,
including functions supported between the nuclei or at infinity.
\end{proposition}

\begin{proof}
Choose smooth real functions \(\theta_0,\theta_1,\ldots,\theta_m\)
with
\[
 \sum_{a=0}^m\theta_a^2=1,\qquad
 \theta_i=1\ \text{on }r_i\le D/5,\qquad
 \supp\theta_i\subseteq\{r_i\le D/4\}.
\]
The balls of radius \(D/4\) are disjoint by
\eqref{eq:continuum-separations}.  In the annulus
\(D/5<r_i<D/4\), take
\(\theta_i=\cos f(r_i/D)\) and
\(\theta_0=\sin f(r_i/D)\), where \(f\) is a fixed smooth
transition from zero to \(\pi/2\), constant near both endpoints.
Set \(\theta_0=0\) in the inner balls and one outside the outer
balls.  At most one transition annulus is present at any point,
so
\begin{equation}
 \sum_{a=0}^m|\nabla\theta_a|^2\le CD^{-2}.
 \label{eq:continuum-ims-gradients}
\end{equation}
This construction also verifies smoothness at the boundaries.

Let \(u\in H^1(\mathbb R^3;\mathbb C^2)\), \(Pu=0\), and
\(\|u\|=1\).  In spin component \(\sigma\), global orthogonality
implies
\[
 b_{i\sigma}:=\langle\psi_i,\theta_i u_\sigma\rangle
 =\langle(\theta_i-1)\psi_i,u_\sigma\rangle.
\]
The tail estimate, integrated over \(r_i\ge D/5\), gives
\[
 \|(\theta_i-1)\psi_i\|
 \le C\poly(D)e^{-.99(.2-.01)D}
 \le e^{-.18D}.
\]
Here the polynomial comes only from radial integration and has
fixed degree.  Consequently
\begin{equation}
 \sum_{i,\sigma}|b_{i\sigma}|^2\le me^{-.36D}.
 \label{eq:continuum-local-ground-leakage}
\end{equation}
There is no assumption that a localized vector is exactly
orthogonal to the local ground state.

On \(\supp\theta_i\), every \(k\ne i\) satisfies
\[
 r_k\ge D-2S-D/4=3D/4-2S>D/8.
\]
Hence \(\chi_k=1\) throughout that support and \(h=h_i\) there.
Choose a fixed
\(0<g_{\mathrm{loc}}\le\min(g_{\mathrm w}/2,1/4)\).
Equations \eqref{eq:continuum-well-gap} and
\eqref{eq:continuum-tail-input}, for large \(D\), imply
\begin{align*}
 \langle\theta_i u,A\theta_i u\rangle
 &\ge g_{\mathrm{loc}}\|\theta_i u\|^2
   -(g_{\mathrm{loc}}+e^{-1.01D})\sum_\sigma|b_{i\sigma}|^2 .
\end{align*}

On \(\supp\theta_0\), each primary distance is at least \(D/5\).
Each secondary distance is at least
\(D/5-\rho_i\ge D/5-4\ge D/6\).
Since \(q+q'\le2\) for large \(D\), the magnitude of the full
attraction on this support is at most \(Cm/D\).  Thus
\[
 \langle\theta_0u,A\theta_0u\rangle
 \ge(1/2-Cm/D)\|\theta_0u\|^2
 \ge g_{\mathrm{loc}}\|\theta_0u\|^2.
\]

The kinetic product rule yields the IMS identity
\[
 \langle u,Au\rangle
 =\sum_{a=0}^m\langle\theta_au,A\theta_au\rangle
 -\frac12\int\sum_{a=0}^m|\nabla\theta_a|^2|u|^2.
\]
It holds first for smooth functions and then on \(H^1\) by
continuity of the Coulomb forms.  Summing the preceding bounds
and using \eqref{eq:continuum-ims-gradients} and
\eqref{eq:continuum-local-ground-leakage} gives
\[
 \langle u,Au\rangle
 \ge g_{\mathrm{loc}}-Cme^{-.36D}-CD^{-2}\ge g_{\mathrm{loc}}/2.
\]
Taking \(g=g_{\mathrm{loc}}/2\) proves the proposition.
\end{proof}

\subsection{Uniform many-fermion bounds and elimination of high modes}

Let
\[
 \mathfrak H_m=\bigwedge\nolimits^m\mathfrak h,\qquad
 \Pi=\bigwedge\nolimits^mP,\qquad Q_m=I-\Pi .
\]
Here \(\Pi\) projects onto \emph{all} \(m\)-electron configurations
in the \(2m\) spin orbitals; it does not impose single occupancy
or a total-spin condition.  Denote second quantization on this
fixed particle-number space by
\(d\Gamma(B)=\sum_{\ell=1}^mB^{(\ell)}\).
In the scaled coordinates, write
\(W_{\mathrm{ee}}=\sum_{1\le\ell<k\le m}|y_\ell-y_k|^{-1}\)
for the electron repulsion, acting trivially on spin.

Under the unitary dilation \(y_\ell=Z_*x_\ell\), the physical
Hamiltonian obeys the exact relation
\begin{equation}
 \frac{H}{Z_*^2}
 =d\Gamma(h)+\frac{W_{\mathrm{ee}}}{Z_*}.
 \label{eq:continuum-exact-dilation}
\end{equation}
We henceforth use this unitary identification and set
\begin{equation}
 \mathcal H=\frac{H}{Z_*^2}+\frac m2
 =d\Gamma(A)+\frac{W_{\mathrm{ee}}}{Z_*}.
 \label{eq:continuum-scaled-hamiltonian}
\end{equation}
To compare its spectral infimum with its compression to \(\Pi\),
we will bound the gap on \(Q_m\) and the coupling
\(Q_m\mathcal H\Pi\). The one-particle estimates control the
contribution from \(d\Gamma(A)\); the following Hardy estimate
controls the repulsion on every vector in \(\operatorname{Ran}\Pi\).

We briefly specify the domains used below.  On the antisymmetric
subspace, the kinetic form domain is \(H^1\) in all \(3m\)
coordinates, with its finite spin components.  Its operator
domain is \(H^2\) in those coordinates.  Translated Hardy
inequalities in one particle coordinate show that each nuclear
or pair Coulomb multiplier is bounded from \(H^1\) to \(L^2\).
For a fixed instance their finite sum therefore satisfies
\(\|Vu\|_2\le C_{\rm inst}\|u\|_{H^1}\).
The Fourier inequality
\(\|\nabla u\|_2\le\varepsilon\|\Delta u\|_2+
C_\varepsilon\|u\|_2\)
makes the potential infinitesimally operator-bounded relative
to the total Laplacian.  The Coulomb operator consequently has
the indicated \(H^2\) operator domain and \(H^1\) form domain.
The same reasoning applies to the one-particle operators.
Every \(\omega_i\) is in \(H^2(\mathbb R^3)\).  Products of these
functions have square-integrable pure second derivatives and
mixed first derivatives, so every vector in \(\operatorname{Ran}\Pi\)
belongs to the many-particle operator domain.  Both \(\Pi\) and
\(Q_m\) preserve the form domain.

\begin{lemma}[A dimension-independent fermionic Hardy bound]
\label{lem:fermion-hardy}
Uniformly over all normalized \(u\in\operatorname{Ran}\Pi\),
\begin{equation}
 \sum_{\ell=1}^m\|\nabla_{y_\ell}u\|_2^2\le Cm,
 \qquad
 \|W_{\mathrm{ee}}\Pi\|\le Cm^2.
 \label{eq:continuum-fermion-hardy}
\end{equation}
The constants do not depend on the dimension
\(\binom{2m}{m}\) of \(\operatorname{Ran}\Pi\).
\end{lemma}

\begin{proof}
Let \((K_{\mathrm{raw}})_{ij}
=\langle\nabla\psi_i,\nabla\psi_j\rangle\).  This is positive
semidefinite, with trace at most \(Cm\).
In the orthonormal spatial modes the kinetic matrix is
\[
 K_0=G^{-1/2}K_{\mathrm{raw}}G^{-1/2}.
\]
Hence
\[
 \Tr K_0
 =\Tr(K_{\mathrm{raw}}G^{-1})
 \le\|G^{-1}\|\Tr K_{\mathrm{raw}}\le Cm.
\]
The spin matrix \(K=K_0\otimes I_2\) has the same bound after
adjusting the constant.

For any normalized fermionic \(u\), its one-particle density
matrix \(\gamma\) in these modes satisfies \(0\le\gamma\le I\).
Indeed, for a coefficient vector \(v\), the quadratic form of
\(\gamma\) is \(\|c(v)u\|^2\); the anticommutation relation
\(c(v)c(v)^\dagger+c(v)^\dagger c(v)=\|v\|^2I\)
places it between zero and \(\|v\|^2\).
Thus, for an arbitrary superposition in the finite-mode space,
\[
 \sum_\ell\|\nabla_{y_\ell}u\|_2^2
 =\Tr(\gamma K)\le\Tr K\le Cm.
\]

For a pair \(\ell<k\), applying the three-dimensional Hardy
inequality in \(y_\ell\), with \(y_k\) and all other variables
held fixed, gives
\[
 \|\,|y_\ell-y_k|^{-1}u\|_2
 \le2\|\nabla_{y_\ell}u\|_2.
\]
Integration over the remaining variables and spin components
justifies this inequality for the full wavefunction.
Put \(b_m=m(m-1)/2\).  Cauchy--Schwarz in the sum of pair
multipliers gives, when \(m\ge2\),
\begin{align*}
 \|W_{\mathrm{ee}}u\|_2^2
 &\le b_m\sum_{\ell<k}
          \|\,|y_\ell-y_k|^{-1}u\|_2^2\\
 &\le4b_m\sum_{\ell<k}\|\nabla_{y_\ell}u\|_2^2\\
 &\le4b_m(m-1)\sum_\ell\|\nabla_{y_\ell}u\|_2^2
 \le Cm^4.
\end{align*}
For \(m=1\) the repulsion is zero.  This proves the operator
bound on the whole finite-mode space, rather than only on its
individual Slater determinants.
\end{proof}

\begin{lemma}[Eliminating the many-particle complement]
\label{lem:continuum-elimination}
For sufficiently large \(D\), let
\(E_\Pi=\min\operatorname{Spec}
(\Pi\mathcal H\Pi|_{\operatorname{Ran}\Pi})\).
Then
\begin{equation}
 0\le E_\Pi-\inf\operatorname{Spec}\mathcal H
 \le e^{-1.2D}.
 \label{eq:continuum-leakage-error}
\end{equation}
More explicitly, the upper bound can be taken to be
\[
 C\left(m^{3/2}e^{-.75D}+m^2/Z_*\right)^2.
\]
\end{lemma}

\begin{proof}
Let \(\varepsilon=\|AP\|\le C\sqrt m\,e^{-.75D}\).
The one-particle off-block operator
\[
 R_{\mathrm{off}}=PA(I-P)+(I-P)AP
\]
extends to a bounded self-adjoint operator, of norm
\(\|(I-P)AP\|\le\varepsilon\).  The adjoint identifies the two
blocks even though \(A\) itself is unbounded.
Its block-diagonal remainder satisfies
\[
 A-R_{\mathrm{off}}
 =PAP+(I-P)A(I-P)\ge-\varepsilon P+g(I-P)
\]
by Proposition~\ref{prop:complement-gap}.  Consequently, on the
full many-particle form domain,
\begin{equation}
 d\Gamma(A)\ge gN_\perp-2m\varepsilon I,
 \qquad N_\perp=d\Gamma(I-P).
 \label{eq:continuum-high-number-bound}
\end{equation}
Here the low diagonal can cost at most \(m\varepsilon\), and
\(\|d\Gamma(R_{\mathrm{off}})\|\le m\varepsilon\).
The commuting one-particle projections \((I-P)^{(\ell)}\)
show that \(N_\perp\) has nonnegative integer eigenvalues,
with kernel exactly \(\operatorname{Ran}\Pi\).  Therefore
\(N_\perp\ge Q_m\), including on superpositions of states
with different numbers of high-mode particles.

The repulsion form is nonnegative on every wavefunction,
so its compression to \(Q_m\) is nonnegative even though
\(W_{\mathrm{ee}}\) need not commute with \(\Pi\).
Equation~\eqref{eq:continuum-high-number-bound} gives
\begin{equation}
 Q_m\mathcal H Q_m\ge(g-2m\varepsilon)Q_m
 \ge cQ_m,\qquad c=g/2,
 \label{eq:continuum-many-body-gap}
\end{equation}
for large \(D\).

By Lemma~\ref{lem:fermion-hardy}, the coupling from the finite
subspace has operator norm at most
\begin{equation}
 \beta:=\|Q_m\mathcal H\Pi\|
 \le m\varepsilon+\frac{Cm^2}{Z_*}
 \le Cm^{3/2}e^{-.75D}+\frac{Cm^2}{Z_*}.
 \label{eq:continuum-coupling}
\end{equation}
The first term follows because the block-diagonal one-particle
sum preserves \(\operatorname{Ran}\Pi\).  The compressed
operator also obeys
\[
 \|\Pi\mathcal H\Pi\|\le m\varepsilon+Cm^2/Z_*=o(1).
\]

For any form-domain vector \(u=p+v\), where \(p=\Pi u\)
and \(v=Q_m u\), the finite-mode operator-domain inclusion and
the bounded coupling justify
\begin{align*}
 \langle u,\mathcal H u\rangle
 &\ge\langle p,\mathcal H p\rangle
       +c\|v\|^2-2\beta\|p\|\|v\|\\
 &\ge\langle p,\mathcal H p\rangle
       -\frac{2\beta^2}{c}\|p\|^2+\frac c2\|v\|^2.
\end{align*}
Equivalently, this is the form inequality
\[
 \mathcal H\ge
 \Pi\mathcal H\Pi-\frac{2\beta^2}{c}\Pi+\frac c2Q_m.
\]
Since the finite block has norm \(o(1)\), the bottom of the
right side is \(E_\Pi-2\beta^2/c\) for large \(D\).
The opposite variational inequality follows by restricting
test vectors to \(\operatorname{Ran}\Pi\).
This proves the explicit error bound.  In the present regime,
with \(Z_*\ge e^D/2\), it is at most \(e^{-1.2D}\).
All lower bounds apply to quadratic forms and hence to spectral
infima; existence of an eigenvector at the bottom of the full
continuum spectrum is not used.
\end{proof}

In particular, the explicit error in this lemma, after
multiplication by \(Z_*\), is bounded by
\begin{equation}
 C\left(
    m^3e^{-.5D}+m^{7/2}e^{-.75D}+m^4e^{-D}
 \right).
 \label{eq:continuum-leakage-final-units}
\end{equation}
The weaker bound \(Z_*e^{-1.2D}\le2e^{-.2D}\) will suffice.
This squaring of the residual against the fixed complement gap
is what permits the final energy amplification.

\subsection{The compressed Coulomb interaction}

The full-space energy is now controlled by its finite-mode compression,
with an error that remains small after multiplication by $Z_*$. To
identify that compression we retain the direct intersite repulsion and
estimate it explicitly. At the present exponentially amplified scales it
has norm $O(m^2/D)$ in Hubbard units, unlike constructions that must keep
it at leading order at polynomial confinement scales.

Identify \(\operatorname{Ran}\Pi\) with
\(\bigwedge^m(\mathbb C^m\otimes\mathbb C^2)\) using the
isometry \(\mathcal J_m=\bigwedge^m\mathcal J\).
Let
\[
 W_{\mathrm{eff}}=\mathcal J_m^*W_{\mathrm{ee}}\mathcal J_m.
\]
Creation and annihilation operators below refer to these
orthonormal abstract site modes.

\begin{lemma}[Coulomb coefficients in the site modes]
\label{lem:continuum-coulomb-compression}
For sufficiently large \(D\),
\begin{equation}
 \left\|W_{\mathrm{eff}}
       -U\sum_{i=1}^m n_{i\uparrow}n_{i\downarrow}\right\|
 \le \frac{Cm^5}{D},\qquad U=\frac58 .
 \label{eq:continuum-coulomb-error}
\end{equation}
\end{lemma}

\begin{proof}
For four spatial functions define
\[
 \mathcal V(a,b;c,d)=
 \iint\frac{\overline{a(x)b(y)}\,c(x)d(y)}
                {|x-y|}\,\dd x\,\dd y.
\]
Hardy in the \(x\) coordinate and Cauchy--Schwarz show that
\begin{equation}
 |\mathcal V(a,b;c,d)|
 \le2\|a\|_2\|b\|_2\|\nabla c\|_2\|d\|_2.
 \label{eq:continuum-coefficient-continuity}
\end{equation}
Telescoping a change of the four arguments therefore makes
this coefficient Lipschitz in their \(H^1\) norms on uniformly
bounded \(H^1\) sets.  The same statement holds for fixed spin
orbitals, with the usual spin Kronecker factors.

The coefficients of \(W_{\mathrm{eff}}\) are these integrals
in the modes \(\omega_i\) and their two spin copies.
Equation~\eqref{eq:continuum-projected-residual} permits their
replacement by \(\phi_i\) with error at most \(Cm/D\) per
coefficient.  Indeed all four \(H^1\) norms stay uniformly
bounded.  In the standard second-quantized expression each
coefficient multiplies
\(\tfrac12c_a^\dagger c_b^\dagger c_d c_c\), a monomial
of norm at most one.  There are at most \((2m)^4\)
coefficients.  Thus this replacement changes the operator
by at most \(Cm^5/D\).
We are comparing coefficient arrays in the fixed abstract
orthonormal site space; no orthogonality of the replacement
functions \(\phi_i\) is asserted or needed.

Choose a fixed smooth radial cutoff equal to one on
\(r_i\le D/5\), zero on \(r_i\ge D/4\), and with gradient
bounded by \(C/D\), and let \(\zeta_i\) be its product with
\(\phi_i\).  Direct integration of the hydrogenic tail gives
\[
 \|\zeta_i-\phi_i\|_{H^1}
 \le C\poly(D)e^{-D/5}\le e^{-.18D}.
\]
The supports of the \(\zeta_i\) are disjoint.  By
\eqref{eq:continuum-coefficient-continuity}, replacing all
\(\phi_i\) by \(\zeta_i\) changes the coefficient operator by
at most \(Cm^4e^{-.18D}\).

For the truncated functions a nonzero coefficient must have
the same site label for both factors at \(x\), and likewise
at \(y\).  Therefore the only surviving terms are onsite
and intersite density terms.  If \(i\ne j\), their spatial
coefficient satisfies
\[
 0\le V_{ij}^{\mathrm{dir}}
 =\iint\frac{|\zeta_i(x)|^2|\zeta_j(y)|^2}{|x-y|}
                  \,\dd x\,\dd y
 \le\frac{1}{D/2-2S}\le\frac CD,
\]
using \(\|\zeta_i\|_2\le1\).
The intersite operator is
\(\sum_{i<j}V_{ij}^{\mathrm{dir}}n_i n_j\),
where \(n_i=n_{i\uparrow}+n_{i\downarrow}\) has norm two.
It consequently has norm at most \(Cm^2/D\).

The onsite spatial integral differs from
\(\mathcal V(\phi_i,\phi_i;\phi_i,\phi_i)=U\)
by at most \(Ce^{-.18D}\).  The two equal-spin monomials
vanish by anticommutation.  The two opposite-spin monomials,
together with the factor \(1/2\) in the interaction, give
exactly \(U n_{i\uparrow}n_{i\downarrow}\).
Combining these estimates yields
\[
 \left\|W_{\mathrm{eff}}
       -U\sum_i n_{i\uparrow}n_{i\downarrow}\right\|
 \le C\left(\frac{m^5}{D}
       +m^4e^{-.18D}+\frac{m^2}{D}\right)
 \le\frac{Cm^5}{D}.
\]
\end{proof}

\Needspace{10\baselineskip}
\subsection{The simulation estimate in physical energy units}

\begin{proposition}[Simulation of the Hubbard spectral minimum]
\label{prop:continuum-simulation}
Let \(E_0=\inf\operatorname{Spec}H\) for the physical molecular
Hamiltonian and let \(E_{\mathrm{Hub}}(t)\) be the minimum of
the half-filled Hubbard Hamiltonian in
Proposition~\ref{prop:hubbard}.  The construction satisfies
\begin{align}
 \left|\frac{E_0}{Z_*}+\frac{mZ_*}{2}
                   -E_{\mathrm{Hub}}(t)\right|
 &\le m\delta_1+\frac{Cm^5}{D}+Z_*e^{-1.2D}
       \label{eq:continuum-simulation-units}\\
 &\le C\biggl[
       m^2\epsilon_{\mathrm{hop}}+\frac{m^5}{D}
       +m e^{-.01D}+m^2e^{-.1D}\notag\\
 &\hspace{35mm}{}+m^2e^{-.6D}+e^{-.2D}
       \biggr].\label{eq:continuum-simulation-final}
\end{align}
In particular, for every prescribed fixed inverse-polynomial
error budget \(\epsilon_*\), the polynomial scale and calibration
precision can be chosen so that the left side is at most
\(\epsilon_*\).
\end{proposition}

\begin{proof}
The exact compression, in the abstract site space, is
\[
 \mathcal J_m^*(Z_*\mathcal H)\mathcal J_m
 =d\Gamma\bigl(Z_*\mathcal J^*A\mathcal J\bigr)
    +W_{\mathrm{eff}}.
\]
The Hubbard hopping term is
\(d\Gamma(\mathsf T\otimes I_2)\).
On \(m\) particles, second quantization of a bounded
one-particle difference increases its norm by at most a
factor \(m\).  Lemma~\ref{lem:one-body-matrix} and
Lemma~\ref{lem:continuum-coulomb-compression} therefore give
\[
 \left\|\mathcal J_m^*(Z_*\mathcal H)\mathcal J_m
             -H_{\mathrm{Hub}}(t)\right\|
 \le m\delta_1+Cm^5/D.
\]
For finite self-adjoint matrices, the difference of their
spectral minima is bounded by the operator norm difference,
by the variational principle.  Lemma~\ref{lem:continuum-elimination}
adds at most \(Z_*e^{-1.2D}\) when the restriction to
\(\operatorname{Ran}\Pi\) is removed.
Finally, \eqref{eq:continuum-scaled-hamiltonian} gives the
exact energy identity
\[
 Z_*\inf\operatorname{Spec}\mathcal H
 =\frac{E_0}{Z_*}+\frac{mZ_*}{2}.
\]
This proves \eqref{eq:continuum-simulation-units}.
Substitute \eqref{eq:continuum-one-body-error} and
\(Z_*\le2e^D\) to obtain
\eqref{eq:continuum-simulation-final}.

The right side involves only fixed polynomial factors in
\(m\), the calibrated hopping error, inverse powers of \(D\),
and the displayed decaying exponentials.
Equation~\eqref{eq:continuum-calibration-interface} and
Proposition~\ref{prop:calibration} allow the hopping error
to be made an arbitrarily small fixed inverse polynomial.
The simultaneous choices, their order, and polynomial bit
cost are specified in Section~\ref{sec:completion}.
No selected-orbital assumption has been imposed on the
spectral infimum in this comparison.
\end{proof}
